% =====================================================================
%  BAB 1: PENDAHULUAN & AKSES SISTEM
% =====================================================================

\chapter{Pendahuluan \& Akses Sistem}

CMS Sekolah adalah sistem manajemen konten (\textit{Content Management System}) berbasis web
yang dibangun di atas platform \textbf{Laravel 11} dan \textbf{Livewire}. Sistem ini dirancang
khusus untuk memenuhi kebutuhan institusi pendidikan dalam mengelola informasi publik sekolah
secara mandiri, aman, dan efisien---tanpa membutuhkan keahlian pemrograman.

Sistem ini mencakup modul-modul seperti Berita, Pengumuman, Agenda, Galeri, Data Kepegawaian,
Manajemen Media, hingga Konfigurasi Website. Seluruh perubahan yang dilakukan melalui panel
admin akan langsung terrefleksi pada halaman publik website sekolah secara \textit{real-time}.

% -------------------------------------------------------------------
\section{Matriks Peran (Role) Pengguna}
% -------------------------------------------------------------------

Akses ke panel admin dibatasi oleh tiga tingkatan peran yang berbeda:

\begin{table}[H]
\centering
\begin{tabularx}{\textwidth}{|l|X|}
\toprule
\textbf{Peran (Role)} & \textbf{Hak Akses} \\
\midrule
\textbf{Super Admin} &
  Akses penuh ke \textbf{seluruh} modul termasuk: Manajemen Pengguna,
  Pengaturan Global, URL Login Kustom, Pods (tipe konten dinamis),
  SEO, dan konfigurasi tampilan. \\
\midrule
\textbf{Admin} &
  Dapat mengelola semua konten (Berita, Pengumuman, Agenda, Galeri,
  Slider, Media, Unduhan, Staf, Fasilitas, Prestasi, Ekstrakurikuler).
  Tidak bisa mengubah Pengaturan inti dan tidak bisa mengelola akun pengguna lain. \\
\midrule
\textbf{Editor} &
  Hanya dapat membuat, mengedit, dan menerbitkan Berita (\textit{Post}).
  Tidak memiliki akses ke modul lain. \\
\bottomrule
\end{tabularx}
\caption{Matriks hak akses pengguna CMS Sekolah}
\label{table:roles}
\end{table}

% -------------------------------------------------------------------
\section{Skenario \& Alur Kerja: Login Pertama Kali}
% -------------------------------------------------------------------

Misalkan Ibu Dewi baru saja ditunjuk sebagai Admin website oleh Kepala Sekolah.
\textbf{Super Admin} (Bapak IT Sekolah) telah membuatkan akun untuknya dengan
\textit{role Admin}.

\begin{enumerate}
  \item Bapak IT memberi tahu Ibu Dewi alamat URL login khusus sekolah, misalnya:
        \texttt{http://situs-sekolah.sch.id/\textbf{portal-admin}}.
  \item Ibu Dewi membuka URL tersebut di peramban dan muncul halaman login.
  \item Ia memasukkan \textbf{Email} dan \textbf{Password} sementara yang diberikan.
  \item Sistem memverifikasi kredensial. Jika salah 5 kali berturut-turut,
        sistem secara otomatis memblokir percobaan login selama 1 menit
        (\textit{rate limiting: 5 percobaan per menit}).
  \item Jika berhasil, Ibu Dewi diarahkan ke halaman \textbf{Dashboard} (\texttt{/admin/dashboard}).
\end{enumerate}

\begin{figure}[H]
  \centering
  \includegraphics[width=0.9\textwidth]{../Img/Backend/screencapture-127-0-0-1-8000-admin-dashboard-2026-06-16-17_51_15.png}
  \caption{Halaman Dashboard Admin --- Ringkasan Statistik Konten}
  \label{fig:dashboard}
\end{figure}

% -------------------------------------------------------------------
\section{Panduan Penggunaan: Login \& Dashboard}
% -------------------------------------------------------------------

\subsection{Proses Login}

\begin{enumerate}
  \item Buka peramban web dan navigasikan ke URL Login yang diberikan oleh Super Admin
        (defaultnya: \texttt{/login}, namun bisa dikustomisasi).
  \item Isi kolom \textbf{Email} dengan alamat surel akun Anda.
  \item Isi kolom \textbf{Password} dengan kata sandi.
  \item Tekan tombol \textbf{Masuk}.
\end{enumerate}

\begin{notebox}[Informasi Keamanan]
  Jika URL Login sudah diubah oleh Super Admin (misalnya menjadi \texttt{/portal-guru}),
  akses melalui \texttt{/login} standar akan menghasilkan halaman \textit{404 Not Found}.
  Pastikan Anda menyimpan tautan login yang benar di bookmark peramban Anda.
\end{notebox}

\subsection{Memahami Dashboard}

Setelah login, Anda akan melihat halaman \textbf{Ringkasan Admin} yang terdiri dari:

\begin{itemize}
  \item \textbf{Kartu Statistik (8 buah)}: Menampilkan total data aktif dari modul:
        \begin{itemize}
          \item \textbf{Berita} --- total artikel yang tersimpan (semua status)
          \item \textbf{Pengumuman} --- total pengumuman
          \item \textbf{Agenda} --- total agenda/kegiatan
          \item \textbf{Galeri} --- total album galeri
          \item \textbf{Program Unggulan} --- total ekstrakurikuler
          \item \textbf{Prestasi} --- total prestasi siswa
          \item \textbf{Staff/Guru} --- total data kepegawaian
          \item \textbf{Media (Gambar)} --- total file gambar di pustaka media
        \end{itemize}
  \item \textbf{Panel Berita Terbaru}: Daftar 5 berita terakhir. Setiap entri menampilkan
        judul, status (\textit{Draft / Published}), dan waktu pembaruan terakhir.
        Terdapat tautan \textbf{Edit} untuk langsung membuka form pengeditan.
  \item \textbf{Panel Pengumuman Terbaru}: Mirip dengan Berita Terbaru, menampilkan
        pengumuman terakhir beserta statusnya.
  \item \textbf{Panel Agenda Mendatang}: Menampilkan agenda yang akan datang,
        diurutkan berdasarkan tanggal \textit{start\_at}.
  \item \textbf{Panel Pengguna Baru}: Menampilkan akun pengguna yang baru ditambahkan
        beserta alamat email mereka (hanya terlihat oleh Super Admin/Admin).
\end{itemize}

% -------------------------------------------------------------------
\section{Logout / Keluar dari Sistem}
% -------------------------------------------------------------------

Untuk keluar dari sesi admin:
\begin{enumerate}
  \item Klik nama pengguna atau ikon profil di sudut kanan atas \textit{sidebar}.
  \item Klik tombol \textbf{Keluar} / \textbf{Logout}.
  \item Sistem akan mengakhiri sesi dan mengarahkan Anda kembali ke halaman Login.
\end{enumerate}

\begin{warnbox}[Keamanan Sesi]
  Jangan pernah meninggalkan sesi admin terbuka di komputer umum atau bersama.
  Selalu \textbf{Logout} setelah selesai bekerja.
\end{warnbox}

% -------------------------------------------------------------------
\section{Penanganan Masalah Umum (Akses \& Login)}
% -------------------------------------------------------------------

\begin{itemize}
  \item \textbf{Halaman Login Menampilkan 404}: Super Admin telah mengubah URL Login
        dari \texttt{/login} menjadi alamat rahasia lain. Hubungi Super Admin untuk
        mendapatkan tautan yang benar.
  \item \textbf{Akun Diblokir / Tidak Bisa Login}: Sistem menerapkan pembatasan
        5 percobaan login per menit. Tunggu 1 menit lalu coba kembali. Jika lupa
        password, hubungi Super Admin untuk melakukan reset.
  \item \textbf{Halaman Dashboard Kosong / Angka Statistik Nol}: Ini normal untuk
        sistem yang baru dipasang. Mulailah mengisi konten melalui modul-modul
        yang tersedia.
\end{itemize}
